\honest{Fighting a Rearguard Action Against the Truth}{Eliezer Yudkowsky}{2008}

In 2000 my younger self began to think technically about how to put a morality into an AI. The reasons for doing so do not matter, except that they show the importance of perfectionism. If you investigate something you may find out about it, and he was now spending his full time on AI morality instead of finding reasons not to.

By 2001 he had still not admitted that his 1997 view was a mistake. He had only improved the old plan by adding Friendly AI as a contingency, for ``the unlikely event that life turns out to be meaningless.'' \nb{\textsc{Creating Friendly AI} (2001) fits this: it calls objective morality ``one of the possibilities.'' The quoted phrase itself could not be traced to a text of the period.} The contingency was a line of retreat. He could doubt his metaethics without it meaning the end of the world. It also let him admit a 10\% chance of having been wrong, then a 20\% chance, instead of the whole mistake at once. \nb{``Update Yourself Incrementally'' (2007) had advised shifting belief ``a little'' on weakly contrary evidence. The post does not say where that advice stops and the rearguard action begins.}

If this sounds too slow, I agree. From 1996 to 2000 his plan had been to get any superintelligence as fast as possible, by every method at once: codelet soup, ad hoc heuristics, evolutionary programming, open source, preferably all together in a Manhattan Project. Adding one more approach could not hurt. His attitude to technology, kept from his childhood technophilia, assumed that any movement toward superintelligence was good without a hint of danger. Once the assumption that AI could not be dangerous was gone, he should have declared a ``Halt, Melt, and Catch Fire,'' taken his weight off every belief built on the wrong assumption, and rethought everything from scratch, as an adult believer must clean house on first noticing that God does not exist. That is ``an art I need to write more about.'' \nb{The author did, in ``Crisis of Faith'', sixteen days later.}

Instead he rehearsed the standard arguments against ``relinquishment''. \nb{\textsc{Creating Friendly AI}: ``I simply do not believe the claim that relinquishment is possible.''} I still think those arguments are ``more or less correct'': it is easier to say what someone is doing wrong than to say what is right, and the government still seems unlikely to do anything sensible about Friendly AI. But I am now more sympathetic to what technophobes say about the downsides of technophilia, and many of his hopes for science and private industry now seem to me equally wrongheaded.

He dropped his 1999 idea of an open-source AI Manhattan Project, but overall the plan stayed: charge in, guns blazing, with his best idea at the time. He never said ``I don't know how to do this,'' or that he needed better knowledge, or that the project was not ready to start coding. The clock was ticking. The Singularity Institute would ``start coding as soon as it's got enough money.'' He still preferred full information-sharing, with secrecy for bad guys and openness for good guys, not knowing that Szilard and Fermi had argued the same question. \nb{They had: Fermi opposed secrecy about fission research, and Szilard won the argument at first.}

My lesson: turning one big ``Oops!'' into a gradual shift in probability leaves no single shock to show that a big change is needed. Each small shift lets you repair the arguments for your strategy while keeping the basic idea. The system absorbs the shock without cracking. In rationality, cracking is good and repair is bad. ``In the art of rationality it's far more efficient to admit one huge mistake, than to admit lots of little mistakes.'' \nb{This is the lesson of ``The Importance of Saying `Oops' '' (2007), which withheld its own case. This post tells the case; here it is the only evidence given for the rule.} There is an instinct, I think, to keep one's plans so as not to thrash around wasting resources, and another to defend what one has argued in public, to avoid the humiliation of being wrong. My younger self was not immune; these impulses only had to nudge his thoughts to do the damage.

Even in 2002 he was not sure the 1997 plan could not have worked; you never know. But a time came when it all fell down. To be continued.
